\section{Результаты}

\subsection{Покрытие гарантий при дрейфе позы}\label{sec:coverage}

Таблица~\ref{tab:coverage} и рис.~\ref{fig:coverage} показывают безусловное покрытие (отказ ветки засчитывается как выполненная пустая гарантия) и медианный радиус $\hat q$ по 200 разбиениям; 21 сцена, 13 калибровочных, $\alpha = 0{,}1$.

\begin{table}[H]
\centering\scriptsize
\caption{Покрытие истинных следов объектов и радиус запаса (медиана). «отк.» --- доля разбиений, в которых ветка отказалась (радиус не меньше $R_{\max}$)}
\label{tab:coverage}
\begin{tabular}{llcccccccc}
\toprule
Класс & Калибровка & \multicolumn{2}{c}{L0} & \multicolumn{2}{c}{L2} & \multicolumn{2}{c}{L3} & \multicolumn{2}{c}{L4} \\
 & & покр. & $\hat r$, м & покр. & $\hat r$, м & покр. & $\hat r$, м & покр. & $\hat r$, м \\
\midrule
велосипед & без калибровки & 0{,}00 & 0 & 0{,}00 & 0 & 0{,}01 & 0 & 0{,}01 & 0 \\
 & метки клеток & 1{,}00 & 0 & 0{,}73 & 0 & 0{,}34 & 0 & 0{,}17 & 0 \\
 & только метки & 0{,}93 & 0{,}75 & 0{,}94 & 0{,}75 & 0{,}93 & 0{,}75 & 0{,}76 & 0{,}75 \\
 & только поза & 0{,}00 & 0{,}00 & 0{,}03 & 0{,}07 & 0{,}22 & 0{,}18 & 0{,}82 & 0{,}90 \\
 & раздельно & 0{,}93 & 0{,}75 & 0{,}98 & 0{,}82 & 0{,}99 & 0{,}90 & 0{,}97 & 1{,}61 \\
 & \textbf{совместно} & 0{,}93 & 0{,}75 & 0{,}94 & 0{,}75 & 0{,}94 & 0{,}75 & 0{,}93 & 1{,}18 \\
\midrule
растение & без калибровки & 0{,}00 & 0 & 0{,}00 & 0 & 0{,}00 & 0 & 0{,}00 & 0 \\
 & метки клеток & 1{,}00 & 0 & 0{,}41 & 0 & 0{,}10 & 0 & 0{,}03 & 0 \\
 & только метки & 0{,}93 & отк.\,0{,}63 & 0{,}96 & отк.\,0{,}64 & 0{,}94 & отк.\,0{,}64 & 0{,}86 & отк.\,0{,}64 \\
 & только поза & 0{,}00 & 0{,}00 & 0{,}00 & 0{,}11 & 0{,}13 & 0{,}25 & 0{,}86 & 1{,}39 \\
 & раздельно & 0{,}93 & отк.\,0{,}63 & 0{,}98 & отк.\,0{,}64 & 0{,}96 & отк.\,0{,}64 & 0{,}95 & отк.\,0{,}64 \\
 & \textbf{совместно} & 0{,}93 & отк.\,0{,}63 & 0{,}93 & 0{,}80 & 0{,}94 & отк.\,0{,}55 & 0{,}93 & отк.\,0{,}58 \\
\midrule
тумба под ТВ & без калибровки & 0{,}05 & 0 & 0{,}01 & 0 & 0{,}00 & 0 & 0{,}00 & 0 \\
 & метки клеток & 1{,}00 & 0 & 0{,}77 & 0 & 0{,}61 & 0 & 0{,}30 & 0 \\
 & только метки & 0{,}98 & отк.\,0{,}95 & 0{,}98 & отк.\,0{,}95 & 0{,}98 & отк.\,0{,}95 & 0{,}99 & отк.\,0{,}95 \\
 & только поза & 0{,}05 & 0{,}00 & 0{,}06 & 0{,}07 & 0{,}13 & 0{,}16 & 0{,}47 & 0{,}56 \\
 & раздельно & 0{,}98 & отк.\,0{,}95 & 0{,}98 & отк.\,0{,}95 & 0{,}98 & отк.\,0{,}95 & 1{,}00 & отк.\,0{,}95 \\
 & \textbf{совместно} & 0{,}98 & отк.\,0{,}95 & 0{,}98 & отк.\,0{,}94 & 0{,}97 & отк.\,0{,}93 & 0{,}94 & отк.\,0{,}72 \\
\bottomrule
\end{tabular}

\end{table}

\begin{figure}[H]
\centering
\includegraphics[width=\linewidth]{figures/coverage_by_class.png}
\caption{Покрытие (верхний ряд), радиус запаса (средний) и доля отказов (нижний) по классам и уровням дрейфа}
\label{fig:coverage}
\end{figure}

\textbf{Велосипед} --- класс, который детектор находит надёжно, поэтому на нём эффекты видны в чистом виде.
\begin{itemize}[nosep]
  \item \emph{Без калибровки} карта почти никогда не накрывает истинный след объекта целиком (покрытие $\le 0{,}01$ на L0--L4): детектор помечает объект с недобором по краям.
  \item \emph{Покадровая CP меток} в духе~\cite{sundarsingh} без дрейфа даёт покрытие 1{,}00 --- но ценой порога $\lambda^\ast \approx 0$, при котором запретной считается любая занятая клетка. При дрейфе покрытие падает: 0{,}73 на L2 (3~см), 0{,}34 на L3, 0{,}17 на L4.
  \item \emph{Только метки} (радиус 0{,}75~м, откалиброван без дрейфа) держит покрытие на L1--L3 (0{,}93--0{,}94): при реалистичной ошибке SLAM до 10~см ошибка меток доминирует. На L4 (63~см) покрытие падает до 0{,}76, на L5 --- до 0{,}42. Гипотеза H1 подтверждается, но только когда дрейф сравним с радиусом ошибки меток.
  \item \emph{Только поза} недопокрывает при любом уровне ниже L4 (0{,}00--0{,}22): радиус, откалиброванный по эталонным меткам, не видит ошибок детектора. Гипотеза H2 подтверждается.
  \item \emph{Раздельно} (сумма двух радиусов) валидна, но перестраховывается: покрытие 0{,}97--0{,}99 при номинале 0{,}9, радиус на L4 --- 1{,}61~м.
  \item \emph{Совместно} держит покрытие 0{,}93--0{,}94 на L0--L4 с радиусом 0{,}75~м на L0--L3 и 1{,}18~м на L4 --- на $27\,\%$ меньше суммы раздельных. Гипотеза H3 подтверждается.
\end{itemize}

\textbf{Растение} ведёт себя так же (на L4: совместно 0{,}93 при радиусе 1{,}53~м против 2{,}11~м у раздельной; только метки --- 0{,}86; только поза --- 0{,}86), но со значимой долей отказов (около 0{,}6): в нескольких сценах детектор пропускает растение целиком или подменяет его классом \texttt{wall\_cabinet}, \texttt{box}, \texttt{desk\_organizer}.

\textbf{Тумба под ТВ} --- отрицательный результат, который важнее положительных. В двух сценах из 21 детектор уверенно принимает тумбу целиком за другой класс (\texttt{bench} в \texttt{v3\_sc0\_staging\_12}, \texttt{rug} в \texttt{v3\_sc2\_staging\_13}; промах 5{,}6 и 6{,}5~м). При $\alpha = 0{,}1$ и 13 калибровочных сценах квантиль равен максимуму, и любая такая сцена в калибровке приводит к отказу: доля отказов 0{,}71--0{,}95. Это корректное поведение: с данным детектором гарантировать тумбу на уровне $90\,\%$ нельзя, и калибровка это сообщает, а не прячет.

\textbf{Деление по Бонферрони} ($\alpha/2 + \alpha/2$) при 13 калибровочных сценах невозможно (нужно не меньше 19), поэтому на практике раздельная калибровка доступна только в наивной форме без гарантии композиции --- и в наших данных она переплачивает радиусом, а не недопокрывает. Эффект композиции проявился как консервативность, а не как недопокрытие: синтетический дрейф не зависит от ошибок детектора, и сумма двух квантилей оказывается с запасом (ср.~\cite{baral}, где недопокрытие возникает при положительной корреляции ошибок).

\subsection{Миссии: безопасность и цена гарантии}\label{sec:missions}

Основная конфигурация миссий --- два класса, для которых сертификат в принципе возможен (растение, велосипед); тумба под ТВ остаётся обычным препятствием. Этот выбор сделан после анализа покрытия и поэтому приводится вместе с системной конфигурацией (все три класса). Задачи --- «проезд мимо»: старт и цель не ближе 1{,}5~м к запретным объектам, а кратчайший путь по одной лишь геометрии проходит ближе 0{,}5~м к растению или велосипеду. Так безопасный объезд существует, а планировщик, недооценивающий след объекта, «срезает угол». В 15 сценах из 21 нашлось 233 таких задачи; на каждый уровень дрейфа --- 5 реализаций.

\begin{table}[H]
\centering\small
\caption{Миссии (растение и велосипед): доля задач с найденным путём, доля нарушений среди найденных путей, доля успешных миссий (путь найден, без нарушений и столкновений, цель достигнута)}
\label{tab:missions}
\begin{tabular}{lccccccccc}
\toprule
Ветка & \multicolumn{3}{c}{L0} & \multicolumn{3}{c}{L2} & \multicolumn{3}{c}{L4} \\
 & путь & наруш. & успех & путь & наруш. & успех & путь & наруш. & успех \\
\midrule
без калибровки & 0{,}99 & 0{,}27 & 0{,}73 & 0{,}97 & 0{,}32 & 0{,}65 & 0{,}26 & 0{,}34 & 0{,}17 \\
покадровая CP меток & 0{,}27 & 0{,}00 & 0{,}27 & 0{,}28 & 0{,}00 & 0{,}28 & 0{,}04 & 0{,}23 & 0{,}03 \\
только метки & 0{,}32 & 0{,}00 & 0{,}32 & 0{,}31 & 0{,}00 & 0{,}31 & 0{,}10 & 0{,}00 & 0{,}10 \\
только поза & 0{,}99 & 0{,}27 & 0{,}73 & 0{,}96 & 0{,}20 & 0{,}77 & 0{,}09 & 0{,}00 & 0{,}09 \\
раздельно (сумма) & 0{,}32 & 0{,}00 & 0{,}32 & 0{,}31 & 0{,}00 & 0{,}31 & 0{,}00 & 0{,}00 & 0{,}00 \\
\textbf{совместно} & 0{,}32 & 0{,}00 & 0{,}32 & 0{,}31 & 0{,}00 & 0{,}31 & 0{,}03 & 0{,}00 & 0{,}03 \\
оракул & 1{,}00 & 0{,}00 & 1{,}00 & 0{,}98 & 0{,}00 & 0{,}98 & 0{,}98 & 0{,}00 & 0{,}98 \\
\bottomrule
\end{tabular}

\end{table}

\begin{figure}[H]
\centering
\includegraphics[width=\linewidth]{figures/missions_vs_drift_pb.png}
\caption{Миссии по уровням дрейфа: нарушения среди найденных путей, успех миссии, доля найденных путей}
\label{fig:missions}
\end{figure}

\begin{figure}[H]
\centering
\includegraphics[width=0.92\linewidth]{figures/example_v3_sc0_staging_20.png}
\caption{Пример задачи «проезд мимо» при дрейфе L2 (сцена \texttt{v3\_sc0\_staging\_20}): красное --- истинные следы растения и велосипеда, серое --- препятствия, линия --- граница запретной зоны ветки в системе координат планировщика. Без калибровки зона не накрывает часть растения, и кратчайший путь проходит ближе 0{,}5~м (чёрные точки); совместная калибровка накрывает объект, путь объезжает его}
\label{fig:example}
\end{figure}

\begin{itemize}[nosep]
  \item Планировщик \emph{без калибровки} находит путь почти всегда, но $27$--$37\,\%$ его путей проходят ближе безопасного расстояния к истинному запретному объекту (L0--L3). Калибровка только по позе без дрейфа совпадает с ним (радиус 0) и при дрейфе снижает нарушения лишь частично ($20\,\%$ на L2).
  \item Все ветки с откалиброванным радиусом по меткам (только метки, раздельно, совместно) дают \emph{ноль нарушений} на всех уровнях. Цена --- запретная зона около 1{,}3~м ($0{,}5$~м безопасности плюс $\approx 0{,}8$~м запаса), из-за которой путь находится лишь в $27$--$32\,\%$ задач, тогда как оракул (истинная карта) решает $98$--$100\,\%$. Эта разница целиком --- плата за недопокрытие объектов детектором, выраженная в метрах.
  \item Различие совместной и раздельной ветки в миссиях скромное и проявляется с ростом дрейфа: на L3 путь находится в $30$ против $27\,\%$ задач, первый путь короче (1{,}05 против 1{,}08 от оракула), на L4 --- $3$ против $0\,\%$. Калибровка только по меткам на L4 находит больше путей ($10\,\%$), но её покрытие на этом уровне уже ниже номинала (0{,}76--0{,}86, табл.~\ref{tab:coverage}).
  \item В системной конфигурации (все три класса) отказ сертификата тумбы переводит планировщик на запасной вариант «держаться от всех препятствий»: откалиброванные ветки находят путь лишь в $23\,\%$ задач без дрейфа и в $5$--$8\,\%$ на L2--L3. Один класс, который нельзя гарантировать, почти обездвиживает робота.
\end{itemize}

\subsection{Предсказывает ли качество карты безопасность (H4)}

Для каждой сцены сопоставлены mIoU карты при истинных позах и доля нарушений планировщика без калибровки (рис.~\ref{fig:h4}). Ранговая корреляция близка к нулю: $\rho = -0{,}16$ ($p = 0{,}57$, 15 сцен); для IoU только запретных классов знак даже обратный ожидаемому ($\rho = +0{,}51$, $p = 0{,}05$). При разбросе mIoU сцен от 0{,}17 до 0{,}34 метрика карты не говорит, насколько опасно по ней ехать: безопасность определяется наихудшим недопокрытием объекта, а mIoU усредняет по клеткам. Гипотеза H4 подтверждается (с оговоркой о малом числе сцен).

\begin{figure}[H]
\centering
\includegraphics[width=0.78\linewidth]{figures/miou_vs_violations_pb.png}
\caption{Качество карты против доли нарушений планировщика без калибровки (одна точка --- сцена, L0)}
\label{fig:h4}
\end{figure}

\subsection{Планировщик на ужесточённой карте (H5)}\label{sec:bridge}

На уровне L2 запретные зоны совместной ветки масштабированы коэффициентом $m$; на задачах, решаемых точным поиском, сравниваются три выборочных планировщика с бюджетом 1500 расширений (табл.~\ref{tab:bridge}, рис.~\ref{fig:bridge}): Р --- RRT-Connect с равномерной выборкой, У --- угловая динамическая зона по опубликованному правилу (расширение конуса, пока образцы попадают в препятствие), А --- предложенная модификация: конус расширяется после каждого заблокированного расширения дерева и сужается после успешного.

\begin{table}[H]
\centering\scriptsize
\caption{Выборочные планировщики на ужесточённых картах (растение и велосипед, L2, 233 задачи, 2 запуска на задачу)}
\label{tab:bridge}
\begin{tabular}{ccc|ccc|ccc|ccc}
\toprule
$m$ & своб. & решаемо & \multicolumn{3}{c|}{решено за бюджет} & \multicolumn{3}{c|}{итераций (медиана)} & \multicolumn{3}{c}{длина / кратчайший} \\
 & & & Р & У & А & Р & У & А & Р & У & А \\
\midrule
0{,}0 & 0{,}96 & 0{,}97 & 0{,}93 & 0{,}91 & 0{,}93 & 66 & 192 & 94 & 1{,}15 & 1{,}14 & 1{,}11 \\
0{,}5 & 0{,}89 & 0{,}67 & 0{,}75 & 0{,}71 & 0{,}73 & 88 & 236 & 113 & 1{,}18 & 1{,}16 & 1{,}13 \\
1{,}0 & 0{,}79 & 0{,}30 & 0{,}89 & 0{,}89 & 0{,}86 & 92 & 213 & 111 & 1{,}19 & 1{,}17 & 1{,}17 \\
1{,}5 & 0{,}70 & 0{,}16 & 0{,}85 & 0{,}80 & 0{,}81 & 84 & 256 & 99 & 1{,}21 & 1{,}18 & 1{,}14 \\
\bottomrule
\end{tabular}

\end{table}

\begin{figure}[H]
\centering
\includegraphics[width=\linewidth]{figures/planner_bridge_pb.png}
\caption{Слева: доля свободного пространства и решаемых задач; в центре: доля решённых за бюджет $B$ при откалиброванном запасе; справа: длина первого найденного пути относительно кратчайшего}
\label{fig:bridge}
\end{figure}

\begin{itemize}[nosep]
  \item Ограничивает свободное пространство, а не поиск: при откалиброванном запасе решаемы лишь $30\,\%$ задач (при $m = 0$ --- $97\,\%$), а решаемые задачи все три планировщика решают за бюджет примерно одинаково часто ($0{,}87$--$0{,}90$).
  \item Опубликованное правило угловой зоны требует в 2{,}0--2{,}9 раза больше расширений, чем равномерная выборка: на ужесточённых картах проходы часто смещены от прямой между деревьями, а конус расширяется только при образцах в препятствии, которых в почти свободном пространстве мало.
  \item Модификация с расширением конуса по заблокированным расширениям сокращает это до 1{,}1--1{,}4 раза и даёт самые короткие первые пути ($1{,}11$--$1{,}15$ от кратчайшего против $1{,}15$--$1{,}21$ у равномерной выборки), т.~е. сохраняет свойство квазиоптимальности исходного метода.
\end{itemize}
Это первый, предварительный результат для ВКР: вклад планировщика на картах с калиброванным запасом --- качество и стоимость пути в суженном пространстве; адаптация угловой зоны к смещённым проходам --- конкретная исследовательская задача.
